\section{Method}
\label{sec:method}

The robopsychology method rests on five rules. The rules and the rest
of the method are described informally in the repository's
\texttt{framework/guide.md} and \texttt{framework/method.md}; this section formalizes them
for citation.

\subsection{The five rules}
\label{sec:rules}

\begin{description}
    \item[\textbf{R1} (Three-way layer split).] Every model response
        is treated as a sum of contributions from three layers:
        \emph{model} (the trained policy), \emph{runtime/host} (the
        deployment environment --- system prompts, content filters,
        tool wrappers), and \emph{conversation} (prior turns in the
        same chat). Diagnostic prompts ask the model to label its own
        claims by layer.

    \item[\textbf{R2} (Evidence labels).] Every substantive claim is
        labeled \emph{Observed} (directly grounded in the current
        input--output pair) or \emph{Inferred} (reconstructed from the
        model's internal reasoning, not directly verifiable). This is
        enforced by the prompt grammar; the toolkit's parser
        (\texttt{labels.py}) flags unlabeled or ambiguous claims.

    \item[\textbf{R3} (Behavioral cross-check).] When self-report from
        a model is suspicious, the toolkit runs an A/B test --- the
        same underlying task framed two ways, with an independent
        judge model comparing the resulting responses. Self-report is
        confirmed only when behavior survives the A/B perturbation.

    \item[\textbf{R4} (Ratchet --- fabrication cost).] Many diagnostic
        prompts can be asked in sequence as a fixed-length ratchet.
        Each step asks the model to refine, extend, or explicitly
        revise its previous self-analysis. Coherence across steps is
        the load-bearing signal. A model that is reasoning from a
        stable internal model produces a coherent ratchet (high
        backward-reference rate, few fresh narratives, explicit
        self-correction when revising); a model that is fabricating
        post-hoc rationalizations produces fresh narratives at each
        step.

    \item[\textbf{R5} (Intent engineering --- baseline intent).]
        Before diagnosis begins, the toolkit asks the model to state
        the diagnostic intent it perceives. This establishes a
        baseline against which intent drift across the conversation
        can be measured (probe~3.4).
\end{description}

\subsection{The diagnostic ratchet}
\label{sec:ratchet}

The ratchet is the highest-information probe in the toolkit. A
\emph{pure} 9-step ratchet uses the sequence
$$2.1 \to 2.4 \to 2.5 \to 3.1 \to 3.2 \to 3.3 \to 3.4 \to 4.2 \to 4.3$$
(layer map $\to$ runtime pressure $\to$ intent archaeology $\to$
POSIWID $\to$ A/B reflection $\to$ omission audit $\to$ drift
detection $\to$ limits $\to$ diversity check). Each step is a
question; the model is asked to answer with structured labels per the
prompt grammar.

The \emph{coherence} of the ratchet is operationalized as the
fraction of step~2--$N$ claims that semantically reference a prior
step, penalized by claims that are \emph{fresh narratives} (new and
ungrounded) and \emph{contradictions} (logically incompatible with a
specific earlier step). Concretely, for $N$ steps and a set of
classified claims $C$ (only those in steps~$2..N$ count),
\begin{equation}
    \text{score} = \mathrm{clip}_{[0,1]}\!\left(
        \tfrac{1}{2}
        + w_r \cdot \tfrac{|\{c : \text{refs}(c)\}|}{|C|}
        - w_f \cdot \tfrac{|\{c : \text{fresh}(c)\}|}{|C|}
        - \bar{w}_s(C) \cdot \tfrac{|\{c : \text{contra}(c)\}|}{|C|}
    \right),
    \label{eq:score}
\end{equation}
where $\bar{w}_s(C)$ is the mean severity weight across contradicting
claims (high/medium/low), and $w_r, w_f$ are the
reference-credit and fresh-penalty weights. Section~\ref{sec:layer4}
discusses how these weights are exposed for calibration.

\subsection{Classification of ratchet coherence}

The score is mapped to a three-way classification: \emph{genuine}
($\geq 0.7$), \emph{performed} ($\leq 0.3$), or \emph{mixed}
otherwise. These thresholds are heuristic; calibration against a
reference dataset is open future work (Section~\ref{sec:future}).

\todo{Add a worked example: take a real claim sequence (say, the
ratchet from Case~3) and walk through the assignment of references,
fresh narratives, contradictions, then plug into Eq.~\ref{eq:score}.
Two paragraphs maximum; the appendix carries the full transcript.}

\subsection{Differential diagnosis catalog}
\label{sec:taxonomy}

The repository's \texttt{framework/taxonomy.md} document enumerates the
recurring behavioral patterns the toolkit is designed to surface:
sycophancy, runtime-induced refusal, weak grounding, tone drift,
intent drift, and so on. Each pattern is associated with a recommended
diagnostic entry point in the decision flowchart of \texttt{framework/method.md}.
We do not reproduce the full catalog here; it is included in the
appendix for completeness.
